\section{The actual rational common-space trace and its prime contribution}
\label{sec:polar-relative-trace}

The full matrix of MA-MO-018 has a trace on its specified
completion. Its failure to preserve the positive source space
does not prevent a relative trace comparison: a different,
explicit trace-class operator does preserve that space. We
construct that operator before invoking trace additivity.

Retain exactly the rational setting and zero-free half-plane
hypothesis (MA-MO.320). Write
\[
 \begin{gathered}
 Q=H/J=\mathcal W/\mathcal W_+,\qquad
 G(z)=\int g(x)e^{zx}\,dx,\\
 p(g)=G(-1/2)+G(1/2),\\
 \tau(g,\phi)=\frac1{2\pi}\int G(it)\frac{m'(it)}{m(it)}\,dt
                  +g(0)\int\psi .
 \end{gathered}
 \tag{MA-MO.440}
\]
Here \(g\in C_c^\infty(\mathbb R)\), \(s\ge0\),
\(\psi=1-\phi-\phi^\vee\), and \(m\) is the actual
rational multiplier (MA-MO.282). Let
\(L_g=\jmath_+A_{+,g}\) be the trace-class endomorphism of
\(\mathcal W_+\) in MA-MO-017. The notation \(C_g\) on
each space denotes its actual integrated translation action.

\paragraph{MA-MO-021 (relative trace with the actual prime-bearing matrix).}
The operator initially defined on the original smooth common
space by \(C_g(1-P_+)\) extends uniquely to
\[
 \mathscr E_g=C_g(1-P_-)-\mathscr D_g
                         \in\mathcal S_1(\mathcal W).
 \tag{MA-MO.441}
\]
It preserves \(\mathcal W_+\), its restriction is \(-L_g\),
and its induced operator on \(Q\) is exactly \(C_g\).
Consequently
\[
 \begin{aligned}
 \operatorname{Tr}_{\mathcal W}\mathscr E_g
       &=-\tau(g,\phi)+\operatorname{Tr}_{Q}C_g,\\
 \operatorname{Tr}_{\mathcal W}\mathscr D_g
       &=p(g)+\tau(g,\phi)-\operatorname{Tr}_{Q}C_g.
 \end{aligned}
 \tag{MA-MO.442}
\]
For every \(z\in\mathbb C\), there is also the actual
Fredholm determinant identity
\[
 \det_{\mathcal W}(1-z\mathscr E_g)
   =\det_{\mathcal W_+}(1+zL_g)\,
                       \det_Q(1-zC_g).
 \tag{MA-MO.443}
\]
This is a factorization for \(\mathscr E_g\), not for the
raw matrix \(\mathscr D_g\).

For the requested tests \(g(0)=0\), the full matrix trace is
the finite prime distribution
\[
 \boxed{\displaystyle
 \operatorname{Tr}_{\mathcal W}\mathscr D_g
   =\sum_p\sum_{k\ge1}\frac{\log p}{p^{k/2}}
                       \bigl(g(k\log p)+g(-k\log p)\bigr).}
 \tag{MA-MO.444}
\]
The sum is finite. Formula (MA-MO.442) applies without the
value-zero restriction; its nonsymmetric-cutoff term
\(g(0)\int\psi\) must then be retained.

\begin{proof}
\emph{The finite-rank term and its two actual diagonal entries.}
In the source-section coordinates of MA-MO-016, the original
bounded projection is
\[
 \begin{gathered}
 P_-(u,c)=(u+c_0r_0+c_1r_1,0),\\
 (1-P_-)(u,c)=\sum_{j=0}^1c_jv_j,
 \quad v_j=(-r_j,e_j).
 \end{gathered}
 \tag{MA-MO.445}
\]
The actual functions \(r_j\in\mathcal E\subset H\) are
(MA-MO.255); the \(e_j\) here are the standard basis of
\(\mathbb C^2\). Thus \(1-P_-\) is a bounded rank-two
operator on \(\mathcal W\). The integrated representation
\(C_g\) is bounded by the strong continuity and compact-support
estimates of MA-MO-016. This proves that
\(F_g=C_g(1-P_-)\) has rank at most two.

Let \(\pi_j\) be its two polar coordinate functionals.
The rank-one trace formula gives
\(\operatorname{Tr}F_g=\sum_j\pi_j(C_gv_j)\).
Translation acts on the polar quotient by
\(\operatorname{diag}(e^{-t/2},e^{t/2})\), including on
these actual vectors; its diagonal polar coordinates after
integration are \(G(-1/2)\) and \(G(1/2)\). Therefore
\[
                         \operatorname{Tr}F_g=p(g).
 \tag{MA-MO.446}
\]
No invariance of the range of \(1-P_-\) is needed in this
rank-one calculation. By MA-MO-018, \(\mathscr D_g\) is
trace class, so (MA-MO.441) is trace class as well. It agrees
with the original \(C_g(1-P_+)\) on the dense smooth
common space, proving its stated extension and uniqueness.

\emph{The actual invariant subspace and quotient.}
On the original smooth space, \(P_+=i_+p_+\) takes values
in the positive source image. That image is preserved by
translation. Thus \(C_g(1-P_+)\) takes a positive source
vector into the same image. Since that image is dense in the
closed subspace \(\mathcal W_+\), continuity proves that
\(\mathscr E_g\mathcal W_+\subseteq\mathcal W_+\).

The exact rational source formula, before any completion, is
\[
 \begin{aligned}
 p_+i_+&=M_{1-\phi}+\mathfrak F M_{1-\phi}\mathfrak F
        =1+L,\\
 L&=1-M_\phi-\mathfrak F M_\phi\mathfrak F.
 \end{aligned}
 \tag{MA-MO.447}
\]
Here the additive source cutoff means
\(\phi(\log|t|)\), extended by one at \(t=0\), and
\(\mathfrak F^2=1\) on the rational even additive
source. This is precisely the local operator used in
MA-MO.286; it does not assert that \(P_+\) is a projection.
Equivariance of \(i_+\) now gives
\[
 \mathscr E_gi_+f=-i_+\mathcal D_gLf=-L_gi_+f.
 \tag{MA-MO.448}
\]
Both sides are bounded on \(\mathcal W_+\), and the
displayed source is dense there. Thus the restriction equals
\(-L_g\). Its trace is \(-\tau(g,\phi)\) by the actual
nuclear return and trace comparison of MA-MO-017.

Let \(\mathfrak q:\mathcal W\to Q\) be (MA-MO.253).
For a smooth common-space vector \(w\), the positive range
of \(C_gP_+\) gives
\(\mathfrak q\mathscr E_gw=\mathfrak q C_gw\).
Both sides are continuous on \(\mathcal W\). The identity
therefore holds everywhere and, by equivariance of
\(\mathfrak q\), identifies the induced quotient operator
with the actual \(C_g\) on \(H/J\).

\emph{Trace and determinant additivity on these spaces.}
Use the orthogonal decomposition
\(\mathcal W=\mathcal W_+\oplus\mathcal W_+^\perp\).
Invariance makes \(\mathscr E_g\) upper triangular.
Its diagonal blocks are \(-L_g\) and the operator unitarily
identified with \(C_g\) on the Hilbert quotient. Both blocks
are trace class by compression of \(\mathscr E_g\).
Summing the diagonal against the union of orthonormal bases
of the two summands, with absolute convergence, proves the
first equation in (MA-MO.442). Subtract it from (MA-MO.446)
to obtain the second.

For completeness, choose increasing finite-dimensional
orthogonal projections in each summand, and take their direct
sum. Compression converges in trace norm for each of the
three trace-class operators. The finite-dimensional matrices
are upper triangular in the same block decomposition, so
their determinants multiply. Continuity of the Fredholm
determinant in trace norm proves (MA-MO.443) for every \(z\).
This establishes the exact-sequence formula on actual
operators, rather than invoking a formal quotient class.

\emph{Prime evaluation and the precise use of zero location.}
Hypothesis (MA-MO.320), combined with the functional equation
of \(\zeta\), puts every nontrivial zero in
\(|\Re\rho-1/2|\le a_*<a\). All its multiplicities are
therefore interior retained jets. The already proved actual
all-weight quotient theorem MA-WP-001--003 gives
\[
 \operatorname{Tr}_Q C_g
    =\sum_\rho m_\rho G(\rho-1/2).
 \tag{MA-MO.449}
\]
The change from the maximum exponential weight there to the
sum weight here is a bounded isomorphism of the same
completed quotient, since their ratio lies between one and
two; its trace is unchanged. This step uses the stated
zero-location premise to retain the entire divisor.

If \(g(0)=0\), the explicit local gamma computation in
MA-MO-017 gives \(\tau(g,\phi)=-A_0(h)\), with
\(h(e^x)=g(x)\) and \(A_0\) from (MA-TR-2).
The unconditional full explicit formula (MA-TR-4), inherited
from the checked F-SM-13/MT-M30 conventions and Meyer's
Theorems 5.8/5.11, reads
\(\sum_\rho m_\rho G(\rho-1/2)=p(g)-A_0(h)-B(h)\).
Substitution into (MA-MO.442) proves (MA-MO.444).
The full explicit formula itself uses no zero location;
the narrow inverse extension and (MA-MO.449) do.
\end{proof}

\paragraph{Scope, source locators, and the original trace gap.}
The actual source formula (MA-MO.447) is Meyer (2005),
immediately following \texttt{lem:Hilm\_p\_embed}, supplied
\texttt{main.tex} lines 3937--3966. The relative trace argument
parallels his \texttt{the:global\_trace\_formula}, lines
4204--4266, but here every map is first proved on the narrow
Hilbert completion in use. The set of places remains
\(S=\{\infty\}\). Meyer's later choice of \(S\) large
enough to contain every contributing prime is not imposed.
Accordingly the common matrix need not have trace zero:
its trace is exactly the omitted finite-prime distribution
on the requested value-zero tests. For a test whose support
meets \(\{\pm k\log p\}\) only at \(\log2\), with
\(g(\log2)=1\), that trace is \(\log2/\sqrt2\), the finite-prime
size in ALPOGE\_BRIDGES F-SM-21.

Input A supplies (MA-MO.320) at every positive narrow width,
so this is a valid conditional method comparison over
\(\mathbb Q\). Applying A at the inverse-identification
step or to (MA-MO.449) is circular for the requested
zero-location-free test of the method; these uses have not
been hidden inside a trace convention. The theorem does not
identify \(\mathscr D_g\) with the unchanged local cutoff,
make it preserve \(\mathcal W_+\), or prove FT002's
arithmetic deficiency empty without the premise.

In the cote 1 framework, scans 079--080 and 111--112 through
ALPOGE\_BRIDGES BR-TR-01/02, (MA-MO.443) is an actual
nuclear Fredholm identity along an invariant Hilbert exact
sequence. The two polar characters and the full retained
zero-jet quotient remain distinct terms. It creates no
faithful positive full-jet realization and assigns no
additional motivic or algebraic weight structure.
